% RESEARCH_DRAFT_SHA256: c93c4716195941dabf94aa2ea53aaae803b364eecf1406a5aab08785e0017596
% Layout and source-reference wrappers expanded for the web renderer.
\documentclass[11pt]{article}
\usepackage[margin=25mm]{geometry}
\usepackage{fontspec,unicode-math}
\setmainfont{Latin Modern Roman}
\setmathfont{Latin Modern Math}
\usepackage{amsmath,amssymb,mathtools,xurl,hyperref}
\hypersetup{colorlinks=true,urlcolor=blue,linkcolor=blue}
\setcounter{secnumdepth}{0}
\setlength{\parindent}{0pt}
\setlength{\parskip}{5pt}
\setlength{\emergencystretch}{3em}
\newcommand{\sref}[2]{\hyperlink{source:#1:#2}{[S#2]}}
\newcommand{\cataloguescope}[1]{\paragraph{Source record.}#1}
\newcommand{\floor}[1]{\left\lfloor#1\right\rfloor}
\begin{document}
\section*{Research-attempt provenance}
This standalone research attempt retains the complete original definition body below,
including its historical source and preparation notes. Those notes describe the
earlier catalogue preparation; the new attempt and its verification are identified
separately after that body. No original ID, catalogue status, or source statement is changed.

\section{Unspecified Flag Algebra Extremal Problem}\label{statement}

\subsection{Definitions and mathematical statement}

\paragraph{Shared notation.}Standard number systems are $\mathbb N,\mathbb Z,\mathbb Q,\mathbb R,\mathbb C$. Specialized notation and index ranges follow the cited source and the definitions below.


A flag-algebra method encodes substructure densities and their consistency relations algebraically, often producing semidefinite relaxations for extremal bounds. A relaxation may contain the original feasible density set and give an upper or lower bound on an extremal quantity.
\paragraph{Record scope and qualifications.}
The catalogue title is retained as metadata. Named certificates, contracts, bounds or frameworks in that title are not asserted as results of the original source extract.
The supplied registry flags this extraction as unrecoverable. It is displayed as a literal source fragment; the definitions provide context rather than a replacement formal statement.

The central quantity ex(C,H), graph conditions, and intended unresolved inequality are not defined by the excerpt.

Source-record limitation: Source flags the statement as unrecoverable: The official report gives a descriptive account of flag-algebra methods but does not formulate a standalone open extremal problem. No separate clean mathematical statement is supplied; this source ID preserves its provenance.
The source text below is rendered literally where necessary; damaged formulas and composite extracts are not silently repaired into a new mathematical claim.
The supplied catalogue marks this record as unrecoverable. Its extract and terminology are documented, but a complete original problem has not been reconstructed.
\paragraph{Statement recorded in the supplied source.}
of convex sets Φk , each containing Φ. Moreover, with appropriate control, opti- mization over the relaxations Φk can be carried out efficiently, for instance via semidefinite programming. This allows us to compute upper bounds for ex(C, H). Many significant results in extremal combinatorics have been obtained via the theory of flag algebras, for instance: computing the minimal number of triangles in graphs with given density, computing the maximum number of pentagons in triangle-free graphs, and obtaining new results regarding the Cacceta-Häggkvist conjecture (see Carli Silva, Oliveira, and Sato [1] for references). Though presented here in the context of graphs, the theory is more general, and can be applied also to digraphs, permutations, and other objects. Conic Linear Optimization for Computer-Assisted Proofs 1071 \href{https://doi.org/10.4171/owr/2022/20}{[S2]}
\paragraph{Source status.}
The supplied catalogue labels this source record open; that metadata was checked against the supplied snapshot on 2026-10-11. The label is preserved for provenance and is not an independent verification of present open status or mathematical correctness. The catalogue literature assessment is dated 2026-08-22.
\subsection{Short English statement}
The fragment describes flag-algebra relaxations and examples of applications.
\subsection{Sources}
\begin{itemize}
\item[S1] Ulam AI and the UnsolvedMath Contributors, UnsolvedMath: A Curated Collection of Open Mathematics Problems, catalogue entry. Catalogue curation and metadata: CC BY 4.0. The source extract is retained; definitions, English summaries and qualifications are editorial additions.
\url{https://huggingface.co/datasets/ulamai/UnsolvedMath}.
\item[S2] Conic Linear Optimization for Computer-Assisted Proofs (2023), Oberwolfach Reports, DOI 10.4171/owr/2022/20, pp. 1039--1089
\url{https://doi.org/10.4171/owr/2022/20}.
\item Source-specific Oberwolfach Reports licences, including CC BY-SA 4.0 where applicable, are retained; see THIRD\_PARTY\_NOTICES.md.
\end{itemize}
\paragraph{Preparation and review.}
The original source and UnsolvedMath contributors retain attribution for the source extract. Codex (GPT-6; exact serving revision not exposed) assisted with definitions, English summaries, scope notes and formatting. No new human mathematical authorship or independently verified solution is asserted. Catalogue identity checked against the complete supplied registry and current repository. Distinct upstream records, including repeated questions and broad or damaged extracts, are retained for complete coverage of the source CS category. A repeated record is not a new independent conjecture. Complete mathematical-equivalence and current-literature review has not been performed. Preparation and final review follow contribution prompts C and D, with source gaps disclosed. Structural and rendering checks do not establish a new theorem or open status.
\label{end:statement}

\section{Literature and context check}

I read the primary report's flag-algebra discussion around pages 1070--1071,
including the relaxation passage from which this record begins mid-sentence.
It describes a method and several applications but specifies no particular
forbidden graph, objective or open extremal constant. The following is a
self-contained auxiliary test case, the classical triangle-free edge bound and
an elementary quantitative stability argument. It is not a repair of the
fragment into an alleged new flag-algebra conjecture.

\section{Research attempt: A finite triangle-free certificate with an explicit stability bound}

\paragraph{An exact local-to-global certificate.}
Let $G$ be a finite simple triangle-free graph with $n\geq1$ vertices and $m$
edges. Write $d_v$ for vertex degrees. For an edge $uv$, the neighborhoods
$N(u)$ and $N(v)$ are disjoint: a vertex in both would make a triangle.
Thus $d_u+d_v\leq n$. Summing over edges counts $d_v$ exactly $d_v$ times,
so
\[
 \sum_v d_v^2\leq nm.
\]
Also $\sum_vd_v=2m$ and
$0\leq\sum_v(d_v-2m/n)^2=\sum_vd_v^2-4m^2/n$.
If $m>0$, the two inequalities imply $m\leq n^2/4$; if $m=0$ this is immediate.
Integrality gives $m\leq\lfloor n^2/4\rfloor$. A complete bipartite graph
with parts of sizes $\lfloor n/2\rfloor,\lceil n/2\rceil$ is triangle-free
and has exactly that many edges. This proves the sharp bound for this
specified auxiliary problem.

\paragraph{Extracting stability rather than only an extremal number.}
Assume $m\geq n^2/4-t$ for some real $t\geq0$. Choose a maximum-degree
vertex $v$, write $a=d_v$, set $A=N(v)$ and $B=V(G)\setminus A$.
The set $A$ is independent. Let $e_B$ count edges inside $B$ and $e_{AB}$
count cross edges. Every vertex in $B$ has degree at most $a$, so
\[
 e_{AB}+2e_B=\sum_{u\in B}d_u\leq a(n-a).
\]
Since $m=e_{AB}+e_B$, it follows that
\[
 m\leq a(n-a)-e_B\leq n^2/4-e_B.
\]
Hence $e_B\leq t$ and $(a-n/2)^2\leq t$.
The number of missing cross edges is
\[
 a(n-a)-e_{AB}=a(n-a)-m+e_B\leq t+e_B\leq2t.
\]
Deleting all $e_B$ internal edges and adding the missing cross edges changes
at most $3t$ edges and produces the complete bipartite graph on $(A,B)$,
whose two part sizes differ by at most $2\sqrt t$.
All of this follows from explicitly summed local inequalities and the displayed
square identity; there is no reliance on a numerical SDP output.

\paragraph{A precise limitation of the route.}
The local inequality $d_u+d_v\leq n$ uses triangle-freeness itself. It cannot
be reused just because some other graph is forbidden: the complete graph on
three vertices has an edge with degree sum $4>3$, while it avoids every
forbidden graph having at least four vertices. Likewise an upper bound obtained
from a relaxation does not automatically yield a stability statement; the
argument above required choosing an actual vertex partition and quantifying
both its internal edges and missing cross edges. No such partition follows
merely from the abstract nesting $\Phi_{k+1}\subseteq\Phi_k$ in the source.

\section{Verification and adversarial audit}

The proof separates $m=0$ before division. The stability argument remains valid
for an edgeless graph: one may take $a=0$, and its hypothesis then forces
$t\geq n^2/4$, making the part-size inequality valid. For odd $n$ the bound
$t=0$ is unattainable, so no false equality assertion is made. The complete
bipartite examples attain the rounded extremal bound.

\paragraph{Audit corrections.} Kept the classical auxiliary target separate from the incomplete source;
proved the stability estimate explicitly rather than inferring it from a numerical
bound, and included empty-edge and parity boundaries.
\paragraph{External computation.} None; symbolic arguments and boundary checks only.
\section{Scope and remaining gaps}

This is a worked finite certificate and its stability extraction for a known
special case. The original damaged passage gives no unique extremal problem
to solve; the argument supplies neither a general flag-algebra hierarchy
convergence theorem nor an answer for the other application examples named in
the fragment. It remains an unresolved source-reconstruction/research record.

\paragraph{Integration review.} Codex (GPT-6; exact serving revision not exposed) reviewed the written arguments and expanded layout and source-reference wrappers. Repository metadata and the production web renderer were checked using Python 3.13 and JavaScript with MathJax 3.2.2. This model-assisted review is not independent human proof verification; no PDF compilation or formal proof checking is asserted.

\section{Completion Estimate}
\noindent\textbf{COMPLETION: 0\%}
The calculations below the preserved source statement investigate a concrete route,
but establish no substantiated new progress toward the entire original target.
Known special cases and failed sufficient conditions are not credited as a solution.
\section{References}
The original source references and catalogue attribution are preserved in the
Sources section above. The following URLs were actually consulted for this attempt;
retention of an additional original reference does not claim fresh consultation of it.
\begin{itemize}
\item \url{https://ems.press/content/serial-article-files/46957} (consulted 2026-10-11). 
\end{itemize}
\section{Model and contributor disclosure}
Codex GPT-6 completed written-proof review of this attempt at its declared scope; this is model-assisted review, not independent human proof checking.
OpenAI gpt-6-astra, requested reasoning configuration ultra, generated this research
attempt on 2026-10-11. The exact serving revision is not exposed. Codex GPT-6
performed orchestration and repository-format review; its exact serving revision
is also not exposed. The mathematical route, boundary checks and final adversarial
self-audit followed contribution prompts A then B in
\texttt{docs/contribution\_prompts.md}, after reading \texttt{CONTRIBUTING.md}.
Python 3.14.2 assembled the standalone documents; any mathematical computation is
specified in Verification. No theorem prover was used. No human mathematical
authorship, substantial human proof correction, or independent human proof checking
is asserted. Model-assisted review, source retrieval, and structural checks do not
establish correctness or novelty. Full private reasoning transcripts are unavailable
as reproducibility artifacts; the complete arguments, inputs and outcomes actually
used are stated here or in the companion batch audit.
\section{Final status}
Unresolved as a new research contribution. This is a concrete route investigation,
not a claim of new complete resolution or a current-status certification. A known
answer reproduced for context does not become a new contribution. The catalogue's
source-status label is preserved; this attempt status does not assert that a
source-labelled solved question is open.
\end{document}
